% AUTHOR: Amlal El Mahrouss
% PURPOSE: WG05.TN001: The Mathematics of Execution.

\documentclass[11pt,a4paper]{article}

\usepackage[utf8]{inputenc}
\usepackage[T1]{fontenc}
\usepackage{amsmath,amssymb,amsthm}
\usepackage{hyperref}
\usepackage[margin=1in]{geometry}

\title{The Mathematics of Execution.}
\author{Amlal El Mahrouss\\
\texttt{amlal@nekernel.org}}
\date{January 2026}

\begin{document}

\bf \maketitle

\begin{center}
	\rule[0.01cm]{17cm}{0.01cm}
\end{center}

\begin{center}
\begin{abstract}
This paper covers the algebraic properties and concepts of `The Execution Semantics' paper by El Mahrouss, A (2026). It is intended as an additional resource over the aforementioned paper.
We will define their concepts alongside their properties. We assume the reader has read the previous paper on 'The Execution Semantics' by El Mahrouss, A.
\end{abstract}
\end{center}

\begin{center}
	\rule[1cm]{17cm}{0.01cm}
\end{center}

\section{The Execution Product}

Let $\Gamma(x, y, z)$ be an execution product of variable arguments $x$, $y$, $z$ defined in compatible and composable execution context $E$: \\
Let the index $\alpha$ denote the current execution domain of an execution product.
Let the following formula:
\begin{equation}
\Gamma(x, y, z) = \prod_{\alpha=1}^{n}(x_{\alpha} \times y_{\alpha} \times z_{\alpha})+C.    
\end{equation}
 Be defined as the execution product. Where $C$ is the Unknown Execution of $\Gamma(x, y, z)$—now defined as $g(E)$.

\subsection{Properties}

\begin{enumerate}

    \item $\Gamma$ as defined previously as the `execution product' shall always be valid within the execution context $E$.
    \item The execution context $E$ shall not denote $\varnothing$.

\end{enumerate}

\section{The Unknown Execution}

Let an `Unknown Execution' $U$ be defined regarding an execution product $\Gamma$. \\
Let $N$ be defined as $\varnothing$.

\subsection{Properties}

\begin{enumerate}

    \item $U$ shall not be equal to $N$.
    \item The execution domain of $\Gamma$ shall not be $N$.

\end{enumerate}

\section{Conclusion}

Such concepts and properties are essential for `Execution Theory''s formality as per defined in El Mahrouss, A. (2026)' paper.

\section{References}

\begin{enumerate}
    \item El Mahrouss, A. (2026). The Execution Semantics: On Axioms, Domains, and Authority. (v2.0.0). Zenodo. https://doi.org/10.5281/zenodo.18366611
\end{enumerate}

\end{document}
